% s-133-p133-m001; original PDF 133
\Aphorism{s-133-p133-m001}{187}
Beim ersten Zuviel des Leidens erfand der Leib die Seele, um nach fernen Tröstungen greifen zu können. Aber nun war Jemand mehr, der litt, und aus einem zweiten Leidens-Übermass müsste die Seele ein Neues erfinden (vielleicht ein dem Traume Ähnliches — oder sollte es die Musik sein?), um mit noch feinerem Geäst an Licht und Luft zu saugen und sich von noch leiseren Winden liebkosen zu lassen. Es ist ihr Zeit gelassen, dieser armen Seele, denn wurzelfest bleibt die Weltesche im Leidensgrunde stehen.
\par

